\begin{definition}
$\GH$ is the category whose objects are globally hyperbolic spacetimes and
whose morphisms $\psi:M\to M'$ are smooth isometric embeddings preserving
orientation and time orientation, with $\psi(M)$ open and causally convex in
$M'$. This is the category $\mathsf{Loc}$ of locally covariant quantum field
theory \citep[\S2]{BFV2003}, \citep[\S2.2]{FV2012}.
\end{definition}

\begin{definition}
A \emph{becoming theory} assigns to each $M\in\GH$ a relation
$B_M\subseteq M\times M$ (read $x\,B_M\,b$ as ``$x$ has become for $b$'') such
that, for every $M$:
\begin{enumerate}
\item[(LC)] for every morphism $\psi:M\to M'$, $x\,B_M\,b\iff
\psi(x)\,B_{M'}\,\psi(b)$;
\item[(T)] $B_M$ is transitive;
\item[(CB)] $\Ichr_M\subseteq B_M$.
\end{enumerate}
\end{definition}

(CB) is Clifton and Hogarth's chronological becoming; they derive it from
``worldline becoming'' and ``worldline transitivity'' \citep[Lemma,
p.~364]{CH1995}, and the derivation carries over verbatim to curved
spacetime. (LC) contains invariance under every orientation- and
time-orientation-preserving isometry of each $M$, which on Minkowski spacetime
$\Mink$ is the proper orthochronous Poincaré group. Its further content is
\emph{locality}: if $N\subseteq M$ is open, causally convex and contains $x$
and $b$, then whether $x$ has become for $b$ is settled by $N$ alone.

\begin{remark}[What (LC) contains]\label{rem:lc}
Three premises should be kept apart. \emph{Naturality}: one rule for all
objects, preserved by every morphism. \emph{Reflection}: the converse
direction of (LC), by which every causally convex region containing $x$ and
$b$ settles the verdict; the proof of Theorem~\ref{thm:main} uses it along $N\hookrightarrow M$ and $N''\hookrightarrow M'$ in part (a), and along $\psi|_N:N\to M'$ and $N''\hookrightarrow\Mink$ in part (b), and functoriality in quantum
field theory supplies no such direction. \emph{Off-shell universality}:
Remark~\ref{rem:offshell}. ``Local'' here is geometric: independence of what
lies around a region, the sense of \citet[\S1]{BFV2003}. It is not
no-signalling or microcausality, and neither of those implies it. The case for
(LC) is that a relation defined from geometry alone should not depend on what
lies outside a causally convex region containing the related events. That is
a premise about becoming, not a consequence of relativity.
\end{remark}

\begin{remark}[The off-shell premise]\label{rem:offshell}
$\GH$ contains every globally hyperbolic metric, not only solutions of
Einstein's equations; the category is defined kinematically in
\citep{BFV2003,FV2012}. Theorem~\ref{thm:main} uses this: its proof builds
auxiliary spacetimes that need not satisfy any field equation with reasonable
matter. This is natural for a kinematic notion of becoming, defined (as Stein
and Clifton--Hogarth define theirs) from geometry alone, but it is a premise.
Section~\ref{sec:energy} examines what happens without it.
\end{remark}
